% Nouveau EGA V (cote 25): archive cote 25, site page 119 (PDF page 120).
% typescript. Draft compared: none.
% AI-transcribed by Claude Opus 5.5 (claude-opus-5-5, adaptive thinking, effort high), Message Batches API, batch ega25b (msgbatch_01A7n5YuXubUPNK1zsepi8RG), 2026-11-13; checked against both supplied scan parts; see CHECK_DECISIONS.tsv.
\scanpage{25/119}
\margin{\unsure{22|2/2}}
\margin{$|$}
Je me répens d'avoir donné la proposition \deleted{\unsure{3|2}} sous une forme éculée et
pisseuse, à mi chemin entre  les hypothèses classiques et les hypothèses
naturelles, et sans donner une récip\unsure{r|p}oque. Aussi je propose l'énoncé:
\uline{Proposition\hand{3}\unsure{\over{2}{3}|3}. \deleted{revue et corrigée}}.  Soit $X \hand{\to} S$ plat loc de présentation
finie,  $Q=\overset{\hand{\vee}}{P} = P(\overset{\hand{\vee}}{E})$  un fibré \xout{vectoriel sur} projectif sur $S$ défini par
un Module loc  libre de type fini $F = \overset{\hand{\vee}}{E}$ ,  considérons l'ensemble
$\hand{\unsure{\phi|\emptyset}}$  des systêmes linéaires \inserted[typed]{$D$ \deleted{\unsure{E}}} de diviseurs \xout{sur} sur $X$, paramétrés par $Q$, tels
que l'ensemble  $Z$ des points fixes de $D$  satisfasse la propriété: $z \hand{\in} Z$
implique prof$_z \underline{O}_{X_s} \hand{\geqslant} 2$  (où  $s$ \xout{\unsure{$= g(z)$)}} est l'image de $z$ dans $S$). \xout{Alors}
\xout{\unsure{$D \hand{\in} \hand{\unsure{\phi|\emptyset}}$ associons le morphisme correspondant $f$ $X$-$Z$ $P$ considérée}}
\xout{\unsure{considérée comme}}  Soit $\hand{\unsure{\psi|\emptyset}}$ l'ensemble des  \deleted{``applications rationelles}\inserted{pseudo-morphismes} rel à
$S$ \deleted{\unsure{*}}\inserted[typed]{$f$} de $X$ dans  $P$, telles que le\xout{\illegible{1}} domaine de définition \inserted{$U=$}\over{dom}{\unsure{U}}(f)
\xout{\unsure{f satisf}} satisfasse la condition : $z \hand{\in} X$-$U$\deleted{(f)} $\hand{\Longrightarrow}$ prof$_z \underline{O}_{X_s} \hand{\geqslant} 2$\inserted{$\unsure{\surd}$, et $f_U = f|U : U \to \check{P}$ satisfait la condition de non dégénérescence envisagée dans 3.1.}. Considé\xout{rons}
rons l'application naturelle  $D \hand{\rightsquigarrow} f_D$  de $\hand{\unsure{\phi|\emptyset}}$  dans $\hand{\unsure{\psi|\emptyset}}$ .  Alors :

a)  Cette application est injective,  et  pour $D \hand{\in} \hand{\unsure{\phi|\emptyset}}$ , l'ensemble des
points fixes $Z$  n'est autre que l'ensemble  complémentaire du domaine de
définition \deleted{\unsure{Dif}} $U$\deleted{(f)} de $f$\hand{$_D$}.

b)  \xout{\unsure{Pour qu'une f $\hand{\unsure{\psi|\emptyset}}$ appartienne à provienne d'une $D$ $\hand{\unsure{\phi|\emptyset}}$, il}} Soit
$f \hand{\in} \hand{\unsure{\psi|\emptyset}}$ , et soit $U$ un ouvert sur lequel $f$ soit défini, et tel que $z \hand{\in} X$-$U$
implique prof $\underline{O}_{X_s,z} \hand{\geqslant} \unsure{2|0}$ \xout{(il en existe} (par exemple $U = U(f)$ domaine de
définition de $f$). Pour que $f$ provienne d'une $D \hand{\in} \hand{\unsure{\phi|\emptyset}}$ , il f et suffit  que\hand{\unsure{$\surd$}}
\deleted{$f_U$ satisfasse la condition énoncée dans 3.1. et que}
posant $\underline{L}_U = f_U^{\hand{\unsure{*|\bullet}}}(\underline{O}_P(1))$  (où  $f_U$: $U \hand{\to} P$  est le morphisme induit par $f$),
\deleted{$f_U$ satisfasse la condition énoncée dans 3.1 et}
\xout{en ait} le Module  $i_{\hand{*}}(\underline{L}_U)$  sur $X$ soit inversible (où $i$:$U \hand{\to} X$ est l'immer-%
sion canonique).

On remarquera que lorsque  les fibres de $X$ sur $S$ satisfont à $(S_2)$,
par exemple sont normales, voire géométriquement  normales, alors la
condition de profondeur  envisagée sur  une partie fermée $Z$ de $X$ dans \deleted{\illegible{1}}
la proposition  signifie simplement que pour tout $s \hand{\in} S$,

\mspage{119}
